\chapter{Cómo Leer Este Libro}
\label{chap:howtoread}

Este libro está organizado en capas: cada parte presupone el conocimiento de
las anteriores. Sin embargo, hemos marcado capítulos autocontenidos para
lectura no lineal.

\section{Notación y Convenciones}

Utilizamos las siguientes convenciones tipográficas:

\begin{itemize}
\item \texttt{\color{keywordblue} código Lisp} -- fragmentos de código
  Common Lisp en el intérprete LisPy.
\item \texttt{\color{keywordblue} código Python} -- fragmentos de código
  Python ejecutables desde LisPy.
\item \textbf{texto en negrita} -- términos nuevos o conceptos clave.
\item \textit{texto en cursiva} -- énfasis o referencias cruzadas.
\end{itemize}

\section{Símbolos de Dificultad}

Cada capítulo está marcado con estrellas que indican su nivel:

\begin{itemize}
\item $\bigstar$\!$\bigstar$\!$\bigstar$\!$\bigstar$ -- \textbf{Principiante.}
  No se requiere experiencia previa en Lisp ni Python.
\item $\bigstar$\!$\bigstar$\!$\bigstar$\!$\bigstar$ -- \textbf{Intermedio.}
  Presupone fundamentos de programación.
\item $\bigstar$\!$\bigstar$\!$\bigstar$\!$\bigstar$ -- \textbf{Avanzado.}
  Requiere comprensión de los conceptos presentados en partes anteriores.
\item $\bigstar$\!$\bigstar$\!$\bigstar$\!$\bigstar$ -- \textbf{Experto.}
  Material avanzado que presupone dominio de las secciones precedentes.
\end{itemize}

\section{Ejecución de Ejemplos}

Los ejemplos de código pueden ejecutarse de tres formas:

\begin{enumerate}
\item \textbf{Línea de comandos:} \texttt{lispy.exe ejemplo.lp}
\item \textbf{REPL interactivo:} \texttt{lispy.exe --repl}
\item \textbf{Evaluación directa:} \texttt{lispy.exe -e "(+ 1 2)"}
\end{enumerate}

\section{Dónde Encontrar el Código Fuente}

El código fuente de todas las librerías discutidas en el libro reside en el
directorio \texttt{lib/} de la distribución de LisPy:

\begin{itemize}
\item \texttt{lib/python-template.lp} -- generación de código Python
\item \texttt{lib/pythonic.lp} -- macro de transformación Lisp a Python
\item \texttt{lib/defpyfun.lp} -- funciones Python con captura Lisp
\item \texttt{lib/numpy-fast.lp} -- aceleración numérica vectorizada
\item \texttt{lib/defnp.lp} -- despacho automático vectorizado/escalar
\item \texttt{lib/sympy-bridge.lp} -- cálculo simbólico con SymPy
\item \texttt{lib/defmath.lp} -- matemáticas híbridas simbólico-numéricas
\item \texttt{lib/persist.lp} -- checkpoint/restore de sesión
\item \texttt{lib/continuations.lp} -- puente de corutinas
\item \texttt{lib/import-hooks.lp} -- transformación de AST
\item \texttt{lib/pandas-bridge.lp} -- DataFrames desde Lisp
\end{itemize}
